\documentclass{ufscar}
\usepackage[utf8]{inputenc}
\usepackage[T1]{fontenc}
\usepackage{geometry}
\usepackage{graphicx}
\usepackage{float}
\usepackage{hyperref}
\usepackage[english,indonesian]{babel}

\title{Optimasi Pemrosesan Template LaTeX Menggunakan Pendekatan Berbasis AI}
\author{Academic Support Agent}
\date{\today}

\begin{document}
\maketitle

\begin{abstract}

\end{abstract}

\chapter{PENDAHULUAN}


\chapter{TINJAUAN PUSTAKA}


% Original content (commented out):
% \documentclass[
	% -- opções da classe memoir --
	12pt,				% tamanho da fonte
	openright,			% capítulos começam em pág ímpar (insere página vazia caso preciso)
	twoside,			% para impressão em recto e verso. Oposto a oneside
	a4paper,			% tamanho do papel. 
	% -- opções da classe abntex2 --
	%chapter=TITLE,		% títulos de capítulos convertidos em letras maiúsculas
	%section=TITLE,		% títulos de seções convertidos em letras maiúsculas
	%subsection=TITLE,	% títulos de subseções convertidos em letras maiúsculas
	%subsubsection=TITLE,% títulos de subsubseções convertidos em letras maiúsculas
	% -- opções do pacote babel --
	english,			% idioma adicional para hifenização
	spanish,			% idioma adicional para hifenização
	brazil				% o último idioma é o principal do documento
	]{ufscar}
	
\usepackage[utf8]{inputenc}
\usepackage{relsize}
\usepackage{lmodern}			% Usa a fonte Latin Modern		
\usepackage[T1]{fontenc}		% Selecao de codigos de fonte.
\usepackage{indentfirst}		% Indenta o primeiro parágrafo de cada seção.
\usepackage{color}				% Controle das cores
\usepackage{graphicx}			% Inclusão de gráficos
\usepackage{microtype} 			% para melhorias de justificação
\usepackage{pdfpages}

% ---
% Pacotes de citações
% ---
\usepackage[brazilian,hyperpageref]{backref}	 % Paginas com as citações na bibl
\usepackage[alf]{abntex2cite}	% Citações padrão ABNT


\usepackage{amsmath}
\usepackage{multirow}
\usepackage[flushleft]{threeparttable}
\usepackage{subcaption}
%\usepackage{graphicx}
\usepackage{lscape}
\usepackage{longtable}
\usepackage{tabularx}
\usepackage{array}
\usepackage[table]{xcolor}
%\usepackage[sectionbib]{bibunits}

% \renewcommand*{\cftchaptername}{CAPÍTULO\space}
 
\definecolor{lightgray2}{RGB}{238,238,238}

% --- 
% CONFIGURAÇÕES DE PACOTES
% --- 

% ---
% Configurações do pacote backref
% Usado sem a opção hyperpageref de backref
\renewcommand{\backrefpagesname}{Citado na(s) página(s):~}
% Texto padrão antes do número das páginas
\renewcommand{\backref}{}
% Define os textos da citação
\renewcommand*{\backrefalt}[4]{
	\ifcase #1 %
		Nenhuma citação no texto.%
	\or
		Citado na página #2.%
	\else
		Citado #1 vezes nas páginas #2.%
	\fi}%
% ---

\newcommand{\textnumero}{
    N\textsuperscript{\underline{o}}
}


% ---
% Informações de dados para CAPA e FOLHA DE ROSTO
% ---
\titulo{TITULO DA DISSERTAÇÃO}
\autor{Fulano de Tal}
\local{Cidade-Estado}
\data{2019}
\orientador[Orientação]{Prof. Dr. Beltrano de Tal}
\instituicao{%
UNIVERSIDADE FEDERAL DE SÃO CARLOS
  \par
  PROGRAMA DE PÓS GRADUAÇÃO EM PLANEJAMENTO E USO DE RECURSOS RENOVÁVEIS}
\tipotrabalho{Dissertação (Mestrado)}
% O preambulo deve conter o tipo do trabalho, o objetivo, 
% o nome da instituição e a área de concentração 
\preambulo{Dissertação apresentada ao Programa de Pós-Graduação em Planejamento e Uso de Recursos Renováveis para obtenção do título de Mestre em Planejamento e Uso de Recursos Renováveis.}
% ---

% ---
% Configurações de aparência do PDF final

% alterando o aspecto da cor azul
\definecolor{blue}{RGB}{41,5,195}

% informações do PDF
\makeatletter
\hypersetup{
     	%pagebackref=true,
		pdftitle={\@title}, 
		pdfauthor={\@author},
    	pdfsubject={\imprimirpreambulo},
	    pdfcreator={LaTeX with abnTeX2},
		pdfkeywords={abnt}{latex}{abntex}{abntex2}{trabalho acadêmico}, 
		colorlinks=true,       		% false: boxed links; true: colored links
    	linkcolor=blue,          	% color of internal links
    	citecolor=blue,        		% color of links to bibliography
    	filecolor=magenta,      		% color of file links
		urlcolor=blue,
		bookmarksdepth=4
}
\makeatother
% --- 

% ---
% Posiciona figuras e tabelas no topo da página quando adicionadas sozinhas
% em um página em branco. Ver https://github.com/abntex/abntex2/issues/170
\makeatletter
\setlength{\@fptop}{5pt} % Set distance from top of page to first float
\makeatother
% ---

% ---
% Possibilita criação de Quadros e Lista de quadros.
% Ver https://github.com/abntex/abntex2/issues/176
%
\newcommand{\quadroname}{Quadro}
\newcommand{\listofquadrosname}{Lista de quadros}

\newfloat[chapter]{quadro}{loq}{\quadroname}
\newlistof{listofquadros}{loq}{\listofquadrosname}
\newlistentry{quadro}{loq}{0}

% configurações para atender às regras da ABNT
\setfloatadjustment{quadro}{\centering}
\counterwithout{quadro}{chapter}
\renewcommand{\cftquadroname}{\quadroname\space} 
\renewcommand*{\cftquadroaftersnum}{\hfill--\hfill}

\addto\captionsbrazil{%
  \renewcommand{\partname}{Capítulo}
}

\setfloatlocations{quadro}{hbtp} % Ver https://github.com/abntex/abntex2/issues/176
% ---

% --- 
% Espaçamentos entre linhas e parágrafos 
% --- 

% O tamanho do parágrafo é dado por:
\setlength{\parindent}{1.3cm}

% Controle do espaçamento entre um parágrafo e outro:
\setlength{\parskip}{0.2cm}  % tente também \onelineskip

% ---
% compila o indice
% ---
\makeindex
% ---

% ----
% Início do documento
% ----
\begin{document}

% Seleciona o idioma do documento (conforme pacotes do babel)
%\selectlanguage{english}
\selectlanguage{brazil}

% Retira espaço extra obsoleto entre as frases.
\frenchspacing 

% ----------------------------------------------------------
% ELEMENTOS PRÉ-TEXTUAIS
% ----------------------------------------------------------
% \pretextual

% ---
% Capa
% ---
\imprimircapa
% ---

% ---
% Folha de rosto
% (o * indica que haverá a ficha bibliográfica)
% ---
\imprimirfolhaderosto*
% ---

% ---
% Inserir a ficha bibliografica
% ---

% Isto é um exemplo de Ficha Catalográfica, ou ``Dados internacionais de
% catalogação-na-publicação''. Você pode utilizar este modelo como referência. 
% Porém, provavelmente a biblioteca da sua universidade lhe fornecerá um PDF
% com a ficha catalográfica definitiva após a defesa do trabalho. Quando estiver
% com o documento, salve-o como PDF no diretório do seu projeto e substitua todo
% o conteúdo de implementação deste arquivo pelo comando abaixo:
%
% \begin{fichacatalografica}
%     \includepdf{fig_ficha_catalografica.pdf}
% \end{fichacatalografica}

\begin{fichacatalografica}

%%%%%%%%%%%%%%%% ATENCAO %%%%%%%%%%%%%%%%%%%%%%%
% Apos gerar ficha catalografica remove comentário a linha abaixo e fazer upload do arquivo ficha_catalografica.pdf %	

%\includepdf{ficha_catalografica.pdf}

\end{fichacatalografica}
% ---

% ---

% ---
% Inserir folha de aprovação
% ---

% Isto é um exemplo de Folha de aprovação, elemento obrigatório da NBR
% 14724/2011 (seção 4.2.1.3). Você pode utilizar este modelo até a aprovação
% do trabalho. Após isso, substitua todo o conteúdo deste arquivo por uma
% imagem da página assinada pela banca com o comando abaixo:
%
% \begin{folhadeaprovacao}
% \includepdf{folhadeaprovacao_final.pdf}
% \end{folhadeaprovacao}
%
\begin{folhadeaprovacao}

%%%%%%%%%%%%%%%% ATENCAO %%%%%%%%%%%%%%%%%%%%%%%
%Remover comentario da linha abaixo e fazer upload do arquivo folha-de-aprovacao.pdf %	

% \includepdf{folha-de-aprovacao.pdf}
  
\end{folhadeaprovacao}
% ---

% ---
% Dedicatória
% ---
\begin{dedicatoria}
   \vspace*{\fill}
   \centering
   \noindent
   \textit{ Dedico este trabalho a todos os alunos de pós-graduação da UFSCar que queiram utilizar o LaTex.} \vspace*{\fill}
\end{dedicatoria}
% ---

% ---
% Agradecimentos
% ---
\begin{agradecimentos}
Agradecimentos

\end{agradecimentos}

%CAPES
%O presente trabalho foi realizado com apoio da Coordenação de Aperfeiçoamento de Pessoal de Nível Superior - Brasil (CAPES) - Código de Financiamento 001

% ---

% ---
% Epígrafe
% ---
\begin{epigrafe}
    \vspace*{\fill}
	\begin{flushright}
		\textit{``Toda a obra de um homem, seja em literatura, música, pintura, arquitetura ou em qualquer outra coisa, é sempre um auto-retrato; e quanto mais ele se tentar esconder, mais o seu caráter se revelará, contra a sua vontade.\\
		(Samuel Butler)}
	\end{flushright}
\end{epigrafe}
% ---

% ---
% RESUMOS
% ---

% resumo em português
\setlength{\absparsep}{18pt} % ajusta o espaçamento dos parágrafos do resumo
\begin{resumo}
Resumo em Português

 \textbf{Palavras-chave}: LaTex, ABNT, UFSCar.
\end{resumo}

% resumo em inglês
\begin{resumo}[Abstract]
 \begin{otherlanguage*}{english}
Abstract in English

   \textbf{Keywords}: LaTEX, ABNT, UFSCar
 \end{otherlanguage*}
\end{resumo}

% ---

% ---
% inserir lista de ilustrações
% ---
\pdfbookmark[0]{\listfigurename}{lof}
\listoffigures*
\cleardoublepage
% ---

% ---
% inserir lista de quadros
% ---
\pdfbookmark[0]{\listofquadrosname}{loq}
\listofquadros*
\cleardoublepage
% ---

% ---
% inserir lista de tabelas
% ---
\pdfbookmark[0]{\listtablename}{lot}
\listoftables*
\cleardoublepage
% ---

% ---
% inserir o sumario
% ---
\pdfbookmark[0]{\contentsname}{toc}
\tableofcontents*
\cleardoublepage
% ---


% ----------------------------------------------------------
% ELEMENTOS TEXTUAIS
% ----------------------------------------------------------
\textual

% ----------------------------------------------------------
% Capitulo 1 - Introdução
% ----------------------------------------------------------

\chapter{Introdução}

%Exemplos

\section{Exemplo Seção}

\subsection{Exemplo Subseção}

Exemplo de citações: 
    Autor, Ano: \cite{weberZhang96}
    Ano: \citeonline{weberZhang96}.
    
\textbf{Texto em Negrito}

\emph{Texto em itálico}

Figura:
\begin{figure}[hbt]
\centering
  \includegraphics[scale=0.8]{figures/sample.jpg}
  \caption{\label{figure:figura1}Legenda figura}
\end{figure}

Referência a um label, exemplo: na Figura \ref{figure:figura1}.

% Nova página
\newpage

Itens:
\begin{itemize}
    \item Item A
    \item Item B
    \item Item C
\end{itemize}

Equação:
\begin{equation}
	\label{eq:equacao1}
	  A = {\pi}r^2
\end{equation}

\begin{equation}
	\label{eq:equacao1}
	  V = \frac{4{\pi}r^3}{3}
\end{equation}

Tabela:
\begin{table}[htp]
    \centering
    \caption{Legenda Tabela}
    \begin{tabular}{c|c|c}
         \textbf{Col 1} & \textbf{Col 2} & \textbf{Col 3}  \\
         \hline
         Linha 1 & Linha 1 & Linha 1  \\ 
         Linha 2 & Linha 2 & Linha 2  \\ 
         Linha 3 & Linha 3 & Linha 3  \\ 
         \hline
    \end{tabular}
    
    \label{tab:tabela1}
\end{table}

% ----------------------------------------------------------


% ----------------------------------------------------------
% Capitulo 2 - Objetivos
% ----------------------------------------------------------
\chapter{Objetivos}

% ----------------------------------------------------------
% Capítulo N - 
% ----------------------------------------------------------
\chapter{Capítulo N}

% ----------------------------------------------------------
% Capítulo N + 1 - Conclusão
% ----------------------------------------------------------
\chapter{Conclusão}


% ----------------------------------------------------------
% Referências bibliográficas
% ----------------------------------------------------------
\bibliography{bibliografia}

% ----------------------------------------------------------
% Glossário
% ----------------------------------------------------------
%
% Consulte o manual da classe abntex2 para orientações sobre o glossário.
%
%\glossary

% ----------------------------------------------------------
% Apêndices
% ----------------------------------------------------------

% ---
% Inicia os apêndices
% ---
\begin{apendicesenv}

% Imprime uma página indicando o início dos apêndices
%\partapendices

% ----------------------------------------------------------
% Apêndice A
% ----------------------------------------------------------
\chapter{Apêndice A}
\label{apendice:a}
% Inclui arquivo apendiceA.tex 
%\input{apendiceA}

% ----------------------------------------------------------
% Apêndice B
% ----------------------------------------------------------
\chapter{Apêndice B}
\label{apendice:b}
% Inclui arquivo apendiceB.tex 
%\input{apendiceB}

\end{apendicesenv}
% ---

%---------------------------------------------------------------------
% INDICE REMISSIVO
%---------------------------------------------------------------------
\phantompart
\printindex
%---------------------------------------------------------------------

\end{document}

\end{document}